\section{INTRODUCTION}
Medical licensing examinations such as the US-MLE represent one of the most demanding knowledge challenges in professional education. Candidates must synthesize clinical reasoning across dozens of specialties under timed, high-stakes conditions. Traditional preparation relies on self-directed study, peer quizzing, and expensive tutoring. These resources are not equally accessible to all candidates. This project presents EMMA: Emergency Medicine Mentoring Agent, a conversational medical study agent built to close that gap. In explanation mode, students pose clinical questions in natural language and receive responses grounded in passages from eighteen standard medical textbooks. In quiz mode, EMMA presents authentic US-MLE-style questions, evaluates answers, and tracks per-specialty performance over time. A collaborative filtering recommender then steers students toward their weakest areas.

EMMA is built on a five-component ML pipeline: a SpaCy NER model for clinical entity extraction, a FAISS vectorstore over 36,723 textbook chunks, a TF-IDF + LinearSVC specialty classifier, a BERTopic clustering model for fine-grained topic routing, and a Surprise-based collaborative filtering recommender. All components feed a locally-run Qwen3-4B-Thinking-2507 LLM. The central empirical question is whether textbook-grounded RAG improves a small local LLM's accuracy on MedQA compared to a no-retrieval baseline. A six-run ablation grid shows the answer is condition-dependent: with the right embedding and LLM pairing, RAG delivers an +11\% point accuracy gain; without those conditions, it can hurt. This nuanced finding is the primary contribution of this report.
